\begin{thebibliography}{99}
% -------------------- ONLINE / HYPERLINK REFERENCES --------------------

\bibitem{ref1}
M. S. Alzboon, “Phishing website detection using machine learning,” Dialnet, 2025. [Online]. Available: \url{https://dialnet.unirioja.es/descarga/articulo/9930098.pdf}

\bibitem{ref8}
“Phishing URL detection using machine learning methods,” ResearchGate. [Online]. Available: \url{https://www.researchgate.net/publication/365790574_Phishing_URL_detection_using_machine_learning_methods}

\bibitem{ref9}
“Phishing URL Detection Using XGBoost,” IJRASET, 2024. [Online]. Available: \url{https://www.ijraset.com/research-paper/phishing-url-detection-using-xgboost}

\bibitem{ref10}
“URL-based phishing detection using machine learning,” IJPREMS, 2025. [Online]. Available: \url{https://www.ijprems.com/ijprems-paper/url-based-phishing-detection-using-machine-learning}

\bibitem{ref11}
“Viable Detection of URL Phishing using Machine Learning,” E3S Conferences, 2023. [Online]. Available: \url{https://www.e3s-conferences.org/articles/e3sconf/pdf/2023/67/e3sconf_icmpc2023_01037.pdf}

\bibitem{ref12}
“Phishing website detection using ML (Project Report),” Sathyabama University. [Online]. Available: \url{https://sist.sathyabama.ac.in/sist_naac/documents/1.3.4/1822-b.e-cse-batchno-287.pdf}

\bibitem{ref13}
“Phishing URL detection using XGBoost and custom feature engineering,” IJRASET. [Online]. Available: \url{https://www.ijraset.com/research-paper/phishing-url-detection-using-xgboost-and-custom-feature-engineering}

\bibitem{ref14}
GitHub Project: “Phishing Website Detection using XGBoost and Random Forest.” [Online]. Available: \url{https://github.com/sudoshivesh/Phishing-Website-Detection-using-XGBoost-and-RFM}

\bibitem{ref15}
MDPI, “Multimodal phishing website detection system,” 2026. [Online]. Available: \url{https://www.mdpi.com/2504-4990/8/1/11}

\bibitem{ref16}
Google, “Safe Browsing API Documentation,” 2023. [Online]. Available: \url{https://developers.google.com/safe-browsing}

\bibitem{ref17}
PhishTank, “Phishing Data Repository,” 2024. [Online]. Available: \url{https://www.phishtank.com}

\bibitem{ref18}
Scikit-learn Developers, “Scikit-learn: Machine Learning in Python,” 2024. [Online]. Available: \url{https://scikit-learn.org}

\bibitem{ref19}
XGBoost Documentation, 2024. [Online]. Available: \url{https://xgboost.readthedocs.io}

\bibitem{ref20}
FastAPI Developers, “FastAPI Documentation,” 2024. [Online]. Available: \url{https://fastapi.tiangolo.com}

\bibitem{ref21}
Docker Inc., “Docker Documentation,” 2024. [Online]. Available: \url{https://docs.docker.com}

\bibitem{ref22}
React Developers, “React Documentation,” 2024. [Online]. Available: \url{https://react.dev}

\bibitem{ref23}
Mozilla Foundation, “MDN Web Docs,” 2024. [Online]. Available: \url{https://developer.mozilla.org}

% -------------------- RESEARCH / NO-LINK REFERENCES --------------------

\bibitem{ref2}
A. Prasad and S. Chandra, “PhiUSIIL phishing URL dataset,” UCI Machine Learning Repository, 2024. doi:10.1016/j.cose.2023.103545

\bibitem{ref3}
J. Ma, L. K. Saul, S. Savage, and G. M. Voelker, “Beyond blacklists: Learning to detect malicious web sites from suspicious URLs,” in \textit{Proc. ACM SIGKDD}, 2009.

\bibitem{ref4}
Y. Zhang, J. Hong, and L. Cranor, “Cantina: A content-based approach to detecting phishing websites,” in \textit{WWW Conference}, 2007.

\bibitem{ref5}
O. K. Sahingoz, E. Buber, O. Demir, and B. Diri, “Machine learning based phishing detection from URLs,” \textit{Expert Systems with Applications}, 2019.

\bibitem{ref6}
D. Sahoo, C. Liu, and S. C. Hoi, “Malicious URL detection using machine learning,” \textit{ACM Computing Surveys}, 2017.

\bibitem{ref7}
T. Chen and C. Guestrin, “XGBoost: A scalable tree boosting system,” in \textit{Proc. ACM SIGKDD}, 2016.

\end{thebibliography}